\documentclass[preprint,12pt]{elsarticle}
\usepackage[T1]{fontenc}
\usepackage{lmodern}
\usepackage{amsmath,amssymb,amsthm}
\usepackage{booktabs}
\usepackage{graphicx}
\usepackage[hidelinks]{hyperref}
\usepackage{xcolor}

\newcommand{\R}{\mathbb{R}}
\newcommand{\bn}{\mathbf{n}}
\newcommand{\bx}{\mathbf{x}}
\newcommand{\by}{\mathbf{y}}
\newcommand{\bu}{\mathbf{u}}
\newcommand{\code}[1]{\texttt{#1}}
\newcommand{\aff}{\operatorname{aff}}
\newtheorem{proposition}{Proposition}
\theoremstyle{definition}
\newtheorem{remark}{Remark}

\journal{Journal of Computational Physics}

\begin{document}

\begin{frontmatter}
\title{Finite volumes on polyhedral cells: exact moments, polynomial reconstruction and the
shallow-water equations}

\author{Oswald Knoth}

\begin{abstract}
We describe the finite-volume method of the Fortran code POLYEDER as it is carried over to the
Julia package \code{dmplex}: cell complexes of general polyhedra with vertex coordinates, the exact
integration of monomials over cells and faces, the geometry derived from these moments, a
polynomial reconstruction that preserves the cell averages, face gradients from interaction regions
around the vertices, and the rotating shallow-water equations on a single layer of prisms over the
sphere. The moments are computed by homogeneous numerical integration: Euler's identity and the
divergence theorem in the affine hull of each cell, face and edge reduce the integral of a monomial
to the vertices, following the cone hierarchy of the cell complex, without the degenerate cases of
vertex formulas and for non-convex cells.
\textcolor{red}{Working draft.}
\end{abstract}

\begin{keyword}
finite volumes \sep polyhedral meshes \sep integration of polynomials \sep $k$-exact reconstruction
\sep multipoint flux approximation \sep shallow-water equations
\end{keyword}
\end{frontmatter}

%==============================================================================
\section{Introduction}\label{sec:intro}

POLYEDER is a finite-volume code for general polyhedral cells: cut cells over orography, Voronoi and
cubed-sphere meshes of a spherical shell, refined and coarsened meshes. Scalars are cell averages,
the momentum is the normal component at the faces. Three ingredients carry the method:
\begin{enumerate}
\item the polynomial reconstruction of cell averages, for the upwind fluxes of the transport of
  scalars and momentum, which needs the moments $\int_C (\bx-\bx_0)^\alpha$ of the cells of a
  stencil;
\item the gradient and the interpolation to the faces from local problems in the interaction
  region around each vertex (the corners of the cells at the vertex), a multipoint flux approximation;
\item the velocity at the cells from the normal components at the faces and the projection of the
  cell tendencies back to the faces.
\end{enumerate}
The shallow-water equations are solved on one layer of prisms over the sphere with free-slip top and
bottom, the same code as for the three-dimensional Euler equations.

In \code{dmplex} the method works on a cell complex (Section~\ref{sec:complex}) that is built from an
extruded surface mesh and later from general three-dimensional meshes. The moments are computed by
homogeneous numerical integration~\cite{Lasserre1998,Chin2015} (Section~\ref{sec:moments}), the
geometry of cells and faces from the moments of degree one (Section~\ref{sec:geometry}).

%==============================================================================
\section{Cell complexes}\label{sec:complex}

A three-dimensional cell complex consists of cells, faces, edges and vertices. The \emph{cone} of a
cell is the list of its faces, that of a face the list of its edges, that of an edge its two
vertices, each with an incidence number $\pm1$ of the cellular boundary,
\begin{equation}
\partial C = \sum_{F\in\mathrm{cone}(C)} s_{CF}\,F,\qquad \partial F = \sum_{E\in\mathrm{cone}(F)}
s_{FE}\,E,\qquad \partial E = v_2 - v_1,
\end{equation}
such that $\partial\partial = 0$. The orientation of a face is the circulation of its oriented
boundary, its normal follows by the right-hand rule, and $s_{CF}=+1$ when the normal points out of
$C$. A face between two cells has opposite signs in their cones, a face on the boundary of the
domain one sign.

\code{dmplex} stores a surface mesh as a \code{Plex} (cells, edges and vertices with cones and
supports, the analogue of PETSc's \code{DMPlex}). Its extrusion \code{ExtrudedPlex} into $n$ layers is
the complex of the tensor prisms: a base cell times a layer is a cell, times a level a horizontal face;
a base edge times a layer a lateral face, times a level a horizontal edge; a base vertex times a layer
a vertical edge, times a level a vertex. Only the surface mesh and $n$ are stored; cones and signs follow
from those of the base, $\partial(p\times I) = \partial p\times I + (-1)^{\dim p}\,p\times\partial I$.
The \code{PolyComplex} holds the cones of the three dimensions in compressed form with the vertex
coordinates; \code{extruded\_complex} places the vertex $v$ of the base at the level of height $z_k$
at $(a+z_k)\,v$ on a sphere of radius $a$ and at $v+z_k\hat{\mathbf z}$ on a plane. A general
three-dimensional mesh fills the same arrays.

%==============================================================================
\section{Exact moments of polytopes}\label{sec:moments}

\subsection{Homogeneous integration}
Let $P\subset\R^3$ be a $k$-dimensional polytope ($k=1,2,3$: a segment, a planar polygon, a
polyhedron), $H=\aff P$ its affine hull and $f$ a homogeneous polynomial of degree $q$, so that
$\by\cdot\nabla f = q f$ (Euler). For a point $\by_0\in H$ the vector field $X=\by-\by_0$ is
tangent to $H$ with $\operatorname{div}_H X = k$, and the divergence theorem in $H$ gives
\begin{equation}
\int_P \operatorname{div}_H(Xf) = \int_P \big(k f + (\by-\by_0)\cdot\nabla f\big)
= (k+q)\int_P f - \int_P \by_0\cdot\nabla f
= \sum_{G} d_G\int_G f ,
\end{equation}
where $G$ runs over the facets of $P$ with the outward unit normals $\nu_G$ in $H$, and
$d_G = (\by-\by_0)\cdot\nu_G$ is constant on $G$ (the distance of $\aff G$ from $\by_0$). For the
monomials $\by^\alpha$, $|\alpha|=q$, this is the recursion
\begin{equation}\label{eq:recursion}
\boxed{\;(k+|\alpha|)\int_P \by^\alpha = \sum_G d_G\int_G \by^\alpha
+ \sum_{i=1}^3 y_{0,i}\,\alpha_i\int_P \by^{\alpha-e_i}\;}
\end{equation}
in the degree for $P$ and down to the vertices in the dimension. With $\by_0$ the foot of the origin
on $H$ the last term vanishes for $k=3$ ($\by_0=0$), and on a segment $[\mathbf a,\mathbf b]$ with the
unit tangent $\mathbf t$ the facets are its end points with $d_{\mathbf b}=\mathbf t\cdot\mathbf b$,
$d_{\mathbf a}=-\mathbf t\cdot\mathbf a$ and $\int_{\{\mathbf v\}}\by^\alpha = \mathbf v^\alpha$.
The order-zero moments are the measures: $k\,|P| = \sum_G d_G |G|$, the formula of the volume of a
pyramid. For a polygon with the vertices $\by_1,\dots,\by_m$ in the order of its orientation, the
normal $\bn = \mathbf A/|\mathbf A|$ follows from the area vector $\mathbf A=\frac12\sum_i
\by_i\times\by_{i+1}$ (Newell), and the outward normal of the edge from $\by_i$ to $\by_{i+1}$ in its
plane is $\nu = \mathbf t\times\bn$.

\begin{remark}
POLYEDER integrates the ridge monomials $(\mathbf b\cdot(\bx-\bx_0))^n$ over a cell by a vertex
formula with weights of the vertices that depend on $\mathbf b$; faces and edges normal to $\mathbf b$
need separate cases (\code{SwitchInt}). The recursion~\eqref{eq:recursion} has no such cases, holds
for non-convex cells and yields all monomials $\by^\alpha$ of a degree at once; a ridge monomial is a
linear combination of them.
\end{remark}

\subsection{Frames and non-planar faces}
The moments of a cell are taken in the local coordinates $\by=(\bx-\bx_C)/h_C$ with $\bx_C$ the mean
of its vertices and $h_C$ the largest distance of a vertex from it, so that $|\by|\le1$ on the cell
and the recursion is well conditioned for any position and size of the cell: the returned moments are
$\int_C \big((\bx-\bx_C)/h_C\big)^\alpha\,dV$. Faces of the cell are integrated in the frame of the
cell. A face whose vertices do not lie in one plane --- a spherical polygon with more than three
vertices --- is replaced by the fan of triangles from the mean of its vertices, oriented as the face;
both cells of the face see the same triangles, so the complex stays closed. The lateral faces of the
prisms are planar, since both of their vertical edges lie on radii.

\subsection{Tests}
The moments of all monomials up to degree four are exact to round-off on the unit cube, on a
triangular prism and on a non-convex L-shaped prism (\code{test/polyhedral.jl}), as is the divergence
theorem $\int_C\partial_i\by^\alpha = \sum_F s_{CF}\, n_{F,i}\int_F\by^\alpha$ with the face moments.

%==============================================================================
\section{Geometry}\label{sec:geometry}

The volume and the centroid of a cell are its moments of degree zero and one, the area vector of a
face is the sum of those of its planar pieces, its centroid follows from the face moments. The
outward area vectors of every cell sum to zero (to $10^{-17}$, Table~\ref{tab:shell}), which makes the
divergence of a constant field vanish exactly. On a spherical shell of thickness $d$ the prisms have
flat faces inscribed in the shell; their total volume converges at second order to that of the shell,
on triangles, hexagons and the squares of the cubed sphere.

\begin{table}[h]
\centering
\caption{Single layer of prisms on the unit sphere, thickness $0.1$: relative error of the total
volume against $\frac{4\pi}{3}(1.1^3-1)$ and the largest $|\sum_F s_{CF}\mathbf A_F|$ of a cell.}
\label{tab:shell}
\begin{tabular}{lrrr}
\toprule
mesh & cells & volume error & closure \\
\midrule
icosahedral, level 2 & 320 & $-3.38\times10^{-2}$ & $2\times10^{-17}$ \\
icosahedral, level 3 & 1280 & $-8.61\times10^{-3}$ & $1\times10^{-17}$ \\
icosahedral, level 4 & 5120 & $-2.16\times10^{-3}$ & $5\times10^{-18}$ \\
hexagonal, level 2 & 162 & $-2.66\times10^{-2}$ & $3\times10^{-17}$ \\
hexagonal, level 3 & 642 & $-6.76\times10^{-3}$ & $1\times10^{-17}$ \\
hexagonal, level 4 & 2562 & $-1.70\times10^{-3}$ & $5\times10^{-18}$ \\
cubed sphere, $4\times4$ & 96 & $-6.51\times10^{-2}$ & $6\times10^{-17}$ \\
cubed sphere, $8\times8$ & 384 & $-1.67\times10^{-2}$ & $2\times10^{-17}$ \\
cubed sphere, $16\times16$ & 1536 & $-4.21\times10^{-3}$ & $9\times10^{-18}$ \\
\bottomrule
\end{tabular}
\end{table}

%==============================================================================
\section{Polynomial reconstruction}\label{sec:reconstruction}

\subsection{Average-preserving least squares}
For every cell $C$ we seek the polynomial $p_C=\sum_k c_k\,\by^{\alpha_k}$ of degree $\le r$ in the
local coordinates $\by = Q_C(\bx-\bx_C)/h_C$ ($\bx_C$ the centroid, $h_C$ the largest distance of a
vertex, $Q_C$ orthogonal) that has the average $\bar u_C$ of the cell and fits the averages of the
other cells of its stencil $S_C$ in the least-squares sense~\cite{Barth1990,OllivierGooch2002},
\begin{equation}
\frac{1}{|C|}\int_C p_C = \bar u_C,\qquad \sum_{C'\in S_C\setminus C}
\Big(\frac{1}{|C'|}\int_{C'} p_C - \bar u_{C'}\Big)^2 \to \min .
\end{equation}
With $a_{C'k} = |C'|^{-1}\int_{C'}\by^{\alpha_k}$ (the moments of Section~\ref{sec:moments} in the frame
of $C$) the constraint eliminates the constant, $c_0 = \bar u_C - \sum_{k\ge1}a_{Ck}c_k$, and the
remaining coefficients solve
\begin{equation}
\sum_{k\ge1}(a_{C'k}-a_{Ck})\,c_k = \bar u_{C'}-\bar u_C,\qquad C'\in S_C\setminus C ,
\end{equation}
by the pseudo-inverse with a relative rank tolerance (POLYEDER's \code{TolRank}). The coefficients are
linear in the averages of the stencil, $c = R_C\,\bar u_{S_C}$, computed once.

\paragraph{Stencils} $S_C$ is $C$ with the cells that share a vertex with it, grown ring by ring
until the reduced matrix has full rank and at least $1.5$ times as many rows as monomials. Inside a
mesh two rings suffice for degree three; at boundaries the stencil grows, for instance in the bottom
and top layers of a column, where two rings see only three layers.

\paragraph{Tangential basis} On one layer of cells over the sphere the stencil does not resolve the
radial structure, and the radial monomials make the least-squares problem singular. The option
\code{basis = :tangential} takes $Q_C$ with the radius at the centroid as its third axis and uses only
the monomials without the third coordinate: a reconstruction in the tangent plane, the moments still
exact over the prisms.

\paragraph{Face values} The average of $p_C$ over a face $F$ of $C$ is
$|F|^{-1}\int_F p_C = \sum_k c_k\, |F|^{-1}\int_F\by^{\alpha_k}$; for the two cells of every face
these are two sparse matrices from the cell averages to the face values (\code{face\_value\_matrices}),
and at the points of a face quadrature two more (\code{point\_value\_matrices}). On a distorted
triangle mesh of a box, extruded into five layers, all face values of polynomials of degree $\le3$ are
reproduced to $10^{-12}$.

\subsection{Transport}
The transport $\partial_t q+\nabla\cdot(q\bu)=0$ of the cell averages takes the upwind values at the
faces. With one volume flux $U_F=\int_F\bu\cdot\bn$ per face, as for the normal velocities of a
staggered scheme, the flux is $U_F\,\bar q_F$ with the upwind face average. This is only second order,
whatever the degree: $\int_F q\,\bu\cdot\bn - U_F\bar q_F = O(h^2)\,|F|$ unless $\bu\cdot\bn$ is
constant on $F$. With the velocity field known at the points $\bx_g$ of a face quadrature
(\code{face\_quadrature}: the seven-point rule of degree five~\cite{Dunavant1985} on the fan of triangles
of the face) the flux $\sum_g w_g(\bu\cdot\bn)_g\,q_g$, upwinded point by point, has the order $r+1$ of
the reconstruction. Time integration is the third-order SSP Runge--Kutta method.

Test fluxes that are exactly divergence free in every cell follow from a vector potential: for the
velocity $\bu = g(\eta)\nabla\psi\times\nabla\eta$ of a stream function $\psi$ of the horizontal position
($\eta$ the radius on the sphere, the height $z$ on a plane) the flux through $F$ is the circulation of
$g(\eta)\psi\nabla\eta$ around $F$, $U_F = \sum_{a\to b}\frac12(\psi(a)+\psi(b))\,(G(\eta_b)-G(\eta_a))$
over the oriented edges of $F$ with $G'=g$; it vanishes on the faces of constant $\eta$, and the fluxes
of a cell sum to zero exactly. On the sphere $\psi=\boldsymbol\Omega\cdot\hat{\mathbf r}$ with
$G(\eta)=\eta^3/3$ gives the solid-body rotation $\boldsymbol\Omega\times\bx$.

\begin{table}[h]
\centering
\caption{Steady flow $\bu=(\sin\pi x\cos\pi y,-\cos\pi x\sin\pi y)$ in the unit square, distorted
triangles ($2n^2$ cells, one layer), $q=\exp(2\psi)$, $\psi=\sin\pi x\sin\pi y/\pi$, relative $L^2$
error at $t=0.5$, fluxes at the quadrature points; CFL $0.4$.}
\label{tab:plane}
\begin{tabular}{rcccccc}
\toprule
 & \multicolumn{2}{c}{degree 1} & \multicolumn{2}{c}{degree 2} & \multicolumn{2}{c}{degree 3} \\
$n$ & error & order & error & order & error & order \\
\midrule
8 & $9.57\times10^{-3}$ & & $4.98\times10^{-3}$ & & $2.71\times10^{-3}$ & \\
16 & $2.09\times10^{-3}$ & 2.20 & $6.97\times10^{-4}$ & 2.84 & $2.16\times10^{-4}$ & 3.65 \\
32 & $4.21\times10^{-4}$ & 2.31 & $9.69\times10^{-5}$ & 2.85 & $1.05\times10^{-5}$ & 4.37 \\
64 & $9.76\times10^{-5}$ & 2.11 & $1.27\times10^{-5}$ & 2.93 & $7.30\times10^{-7}$ & 3.84 \\
\bottomrule
\end{tabular}
\end{table}

With the face fluxes $U_F\bar q_F$ the same flow converges at second order for all degrees ($1.2\times10^{-4}$,
$1.0\times10^{-4}$, $1.3\times10^{-4}$ at $n=64$). On one layer over the sphere (thickness $0.01$) a
Gaussian hill $\exp(-|\hat{\mathbf r}-\hat{\mathbf r}_0|^2/0.4^2)$ is rotated once about an axis tilted by
$45^\circ$ with the divergence-free face fluxes (Table~\ref{tab:sphere}); mass is conserved to $10^{-14}$.
Besides the flux, the flat faces of the prisms inscribed in the shell limit the order to two.

\begin{table}[h]
\centering
\caption{Solid-body rotation of a Gaussian hill over one revolution on one layer of prisms over the
unit sphere, tangential basis, face fluxes $U_F\bar q_F$: relative $L^2$ error.}
\label{tab:sphere}
\begin{tabular}{lrcccc}
\toprule
mesh & cells & degree 0 & degree 1 & degree 2 & degree 3 \\
\midrule
hexagonal 3 & 642 & $7.6\times10^{-1}$ & $3.3\times10^{-1}$ & $3.7\times10^{-1}$ & $1.3\times10^{-1}$ \\
hexagonal 4 & 2562 & $6.3\times10^{-1}$ & $1.2\times10^{-1}$ & $1.2\times10^{-1}$ & $2.2\times10^{-2}$ \\
hexagonal 5 & 10242 & $4.7\times10^{-1}$ & $3.3\times10^{-2}$ & $2.2\times10^{-2}$ & $4.5\times10^{-3}$ \\
icosahedral 3 & 1280 & $6.1\times10^{-1}$ & $2.3\times10^{-1}$ & $1.4\times10^{-1}$ & $9.8\times10^{-2}$ \\
icosahedral 4 & 5120 & $4.6\times10^{-1}$ & $7.0\times10^{-2}$ & $3.1\times10^{-2}$ & $2.0\times10^{-2}$ \\
icosahedral 5 & 20480 & $3.1\times10^{-1}$ & $1.8\times10^{-2}$ & $5.7\times10^{-3}$ & $4.6\times10^{-3}$ \\
\bottomrule
\end{tabular}
\end{table}

\begin{remark}
For the shallow-water equations with the normal momentum at the faces, the mass flux through a face is
a single number, and the transport is second order like that of every staggered scheme; the higher
reconstruction lowers the error constant and the dispersion (Table~\ref{tab:sphere}). Higher order needs
the normal velocity along the face (from the reconstructed cell velocity, or several unknowns per face)
and cells with curved faces on the sphere.
\end{remark}

%==============================================================================
\section{Interaction regions and the face gradient}\label{sec:mpfa}

Around every vertex $v$ the cells at $v$ (the \emph{corners}) and the faces at $v$ form the
interaction region of a multipoint flux approximation~\cite{Aavatsmark2002} (POLYEDER: \code{TGrad}).
Every corner has a linear gradient from the values at the centroids $\bx_F$ of its faces at $v$,
\begin{equation}
\mathbf g_C = G_C\,(u_F-\bar u_C)_{F\ni v},\qquad G_C = \big[(\bx_F-\bx_C)^T\big]_F^{+},
\end{equation}
three faces for the corners of prisms (the pseudo-inverse for more). The values $u_F$ at the faces of
the region are the unknowns of a local linear system: the normal gradient is continuous,
$\bn_F\cdot\mathbf g_{C_1}=\bn_F\cdot\mathbf g_{C_2}$, across every interior face, and vanishes on the
boundary. Its solution expresses the half-face values and normal gradients in the cell averages of
the region; the face value and the normal gradient at a face are their means over its vertices,
weighted by the share of the vertex in the area of the face (the quadrilateral of the vertex, the
midpoints of its edges and the mean of the vertices). The two sparse matrices are exact for linear
fields at all faces whose vertices lie inside the domain (to $10^{-14}$ on distorted triangular prisms
in four layers). On one layer over the sphere the zero normal gradient at the top and the bottom is
exact for every function of the horizontal position.

%==============================================================================
\section{The shallow-water equations}\label{sec:sw}

On one layer of prisms over the sphere, free slip at the top and the bottom, the height $h$ is a cell
average and the normal momentum $U_F=(h\bu)\cdot\bn_F$ lives on the lateral faces:
\begin{align}
\frac{d\bar h_C}{dt} &= -\frac{1}{|C|}\sum_F s_{CF}|F|\,U_F,\\
\frac{dU_F}{dt} &= \bn_F\cdot\big[\mathbf T_C\big]_F - \partial_n\big(\tfrac g2 h^2\big)_F - g\,h_F\,\partial_n b_F,
\qquad
\mathbf T_C = -\frac{1}{|C|}\sum_F s_{CF}|F|\,\frac{U_F}{h_F}\,(\mathbf m)_F^{\mathrm{up}} - f\,\hat{\mathbf r}\times\mathbf m_C .
\end{align}
The pressure gradient, the bathymetry and the face values $h_F$ come from the interaction regions;
the cell momentum $\mathbf m_C$ is reconstructed from the normal components, its Cartesian
components are transported with the upwind values of the polynomial reconstruction
(Section~\ref{sec:reconstruction}, tangential basis), and the cell tendency $\mathbf T_C$ is carried
to the faces with the face values of the interaction regions, again exact for linear fields. Mass is
conserved to round-off. Time integration: the third-order SSP Runge--Kutta method at the gravity-wave
CFL number $0.4$.

\paragraph{The cell momentum} Perot's reconstruction~\cite{Perot2000} (POLYEDER's \code{SumNormal}),
\begin{equation}
\mathbf m_C = \frac{1}{|C|}\sum_F s_{CF}|F|\,U_F\,(\bx_F-\bx_C),
\end{equation}
is exact for constant fields; for linear fields the face-centroid rule misses a quadratic term, an
$O(h)$ error that cancels on symmetric cells only. On triangles the scheme is therefore first order.
A least-squares fit of a linear vector field to the normal components on the faces of the cell and of
its vertex neighbours is exact for linear fields, but unstable in the shallow-water scheme; weighting
the faces of the cell itself twenty times stronger than those of the neighbours makes it stable and
second order on triangles. On hexagons and quadrilaterals Perot's form is more accurate and stays
stable over 15 days, where the weighted fit grows; the default takes the fit on triangles only.

\begin{table}[h]
\centering
\caption{Williamson's case 2 after 5 days on one layer of prisms (thickness 1000 m), reconstruction of
degree 2: relative $L^2$ error of $h$; mass conserved to $10^{-13}$.}
\label{tab:w2}
\begin{tabular}{lrccc}
\toprule
mesh & cells & Perot & weighted fit & default \\
\midrule
icosahedral 3 & 1280 & $3.96\times10^{-3}$ & $2.10\times10^{-3}$ & fit \\
icosahedral 4 & 5120 & $2.00\times10^{-3}$ & $5.68\times10^{-4}$ & \\
icosahedral 5 & 20480 & $1.04\times10^{-3}$ & $1.43\times10^{-4}$ & \\
hexagonal 3 & 642 & $2.74\times10^{-3}$ & $3.47\times10^{-3}$ & Perot \\
hexagonal 4 & 2562 & $6.69\times10^{-4}$ & $1.09\times10^{-3}$ & \\
hexagonal 5 & 10242 & $1.73\times10^{-4}$ & $2.96\times10^{-4}$ & \\
cubed sphere $8^2$ & 384 & $6.63\times10^{-3}$ & $1.03\times10^{-2}$ & Perot \\
cubed sphere $16^2$ & 1536 & $7.67\times10^{-4}$ & $1.22\times10^{-3}$ & \\
cubed sphere $32^2$ & 6144 & $1.84\times10^{-4}$ & $2.85\times10^{-4}$ & \\
\bottomrule
\end{tabular}
\end{table}

\textcolor{red}{Next:} the barotropic instability of Galewsky et al.~\cite{Galewsky2004}, bathymetry
(Williamson's case 5), split-explicit time integration~\cite{Wensch2009}, the GPU, and general
three-dimensional complexes.

\section*{Code availability}
\code{src/Polyhedral/} of \code{dmplex}: \code{Complex.jl} (\code{PolyComplex},
\code{extruded\_complex}), \code{Integration.jl} (\code{MonomialBasis}, \code{segment\_moments},
\code{polygon\_moments}, \code{face\_moments}, \code{cell\_moments}), \code{Geometry.jl}
(\code{PolyGeometry}), \code{Reconstruction.jl} (\code{PolyReconstruction}, \code{face\_value\_matrices},
\code{point\_value\_matrices}), \code{Advection.jl} (\code{PolyAdvection}, \code{stream\_function\_flux}),
\code{InteractionRegion.jl} (\code{interaction\_operators}, \code{cell\_vectors},
\code{linear\_vector\_reconstruction}), \code{ShallowWater.jl} (\code{PolyShallowWater});
tests in \code{test/polyhedral.jl}.

\begin{thebibliography}{99}
\bibitem{Aavatsmark2002} I.~Aavatsmark, An introduction to multipoint flux approximations for
  quadrilateral grids, \emph{Comput. Geosci.} 6 (2002) 405--432.
\bibitem{Barth1990} T.~J. Barth and P.~O. Frederickson, Higher order solution of the Euler equations on
  unstructured grids using quadratic reconstruction, AIAA paper 90-0013 (1990).
\bibitem{Chin2015} E.~B. Chin, J.~B. Lasserre and N.~Sukumar, Numerical integration of homogeneous
  functions on convex and nonconvex polygons and polyhedra, \emph{Comput. Mech.} 56 (2015) 967--981.
\bibitem{Dunavant1985} D.~A. Dunavant, High degree efficient symmetrical Gaussian quadrature rules for the
  triangle, \emph{Int. J. Numer. Methods Eng.} 21 (1985) 1129--1148.
\bibitem{Galewsky2004} J.~Galewsky, R.~K. Scott and L.~M. Polvani, An initial-value problem for testing
  numerical models of the global shallow-water equations, \emph{Tellus} 56A (2004) 429--440.
\bibitem{Lasserre1998} J.~B. Lasserre, Integration on a convex polytope, \emph{Proc. Amer. Math. Soc.}
  126 (1998) 2433--2441.
\bibitem{OllivierGooch2002} C.~Ollivier-Gooch and M.~Van Altena, A high-order-accurate unstructured mesh
  finite-volume scheme for the advection--diffusion equation, \emph{J. Comput. Phys.} 181 (2002) 729--752.
\bibitem{Perot2000} B.~Perot, Conservation properties of unstructured staggered mesh schemes,
  \emph{J. Comput. Phys.} 159 (2000) 58--89.
\bibitem{Wensch2009} J.~Wensch, O.~Knoth and A.~Galant, Multirate infinitesimal step methods for
  atmospheric flow simulation, \emph{BIT Numer. Math.} 49 (2009) 449--473.
\bibitem{Williamson1992} D.~L. Williamson, J.~B. Drake, J.~J. Hack, R.~Jakob and P.~N. Swarztrauber, A standard
  test set for numerical approximations to the shallow water equations in spherical geometry,
  \emph{J. Comput. Phys.} 102 (1992) 211--224.
\end{thebibliography}

\end{document}
